\begin{theorem}[深层规范体积链的 \texorpdfstring{$p$}{p}--进进位重整化]\label{thm:conclusion-foldbin-deep-volume-carry-renormalization}
设 \(k\ge 1\)。记
\[
re_{m+k\to m}:X_{m+k}\to X_m
\]
为 \(k\) 步 restriction，\(d_n(x):=\#\Fold_n^{-1}(x)\)，并定义
\[
\left|G_m^{(2^k)}\right|:=\prod_{w\in X_m}(2^k d_m(w))!,
\qquad
\left|G_{m+k}\right|:=\prod_{u\in X_{m+k}}d_{m+k}(u)!.
\]
则有严格整数恒等式
\[
\frac{\left|G_m^{(2^k)}\right|}{\left|G_{m+k}\right|}
=
\prod_{w\in X_m}
\frac{(2^k d_m(w))!}{\prod_{u\in re_{m+k\to m}^{-1}(w)} d_{m+k}(u)!}
\in \NN,
\]
其中每一因子等于相应后裔多重度束的多项式系数。并且对任意素数 \(p\)，
\[
v_p\!\left(\left|G_m^{(2^k)}\right|\right)-v_p\!\left(\left|G_{m+k}\right|\right)
=
\sum_{w\in X_m}
\mathrm{Carr}_p\!\Bigl(\bigl(d_{m+k}(u)\bigr)_{u\in re_{m+k\to m}^{-1}(w)}\Bigr),
\]
其中 \(\mathrm{Carr}_p\) 表示把该组整数在 \(p\)-进制逐位相加时产生的总进位次数。
\end{theorem}
\begin{proof}
正文已经在一步 restriction 情形证明：\(\left|G_m^{(2)}\right|/\left|G_{m+1}\right|\) 是逐父节点的多项式系数积，而其 \(p\)-进赋值恰等于相应子纤维多重度在 \(p\)-进制相加时的总进位数。

对一般 \(k\)，由 fold--restriction 相容性沿塔复合可知：对每个 \(w\in X_m\)，全部 \(k\) 层后裔满足
\[
\sum_{u\in re_{m+k\to m}^{-1}(w)} d_{m+k}(u)=2^k d_m(w),
\]
因为 \(\Fold_m^{-1}(w)\subset \Omega_m\) 中的每个微态在 \(\Omega_{m+k}\) 中恰有 \(2^k\) 个延拓。于是按父节点分组，
\[
\frac{(2^k d_m(w))!}{\prod_{u\in re_{m+k\to m}^{-1}(w)} d_{m+k}(u)!}
\]
正是一个多项式系数；乘遍全部 \(w\) 即得整数恒等式。

再由 Kummer--Legendre 型多项式系数赋值公式，对每个 \(w\) 都有
\[
v_p\!\left(
\frac{(2^k d_m(w))!}{\prod_{u\in re_{m+k\to m}^{-1}(w)} d_{m+k}(u)!}
\right)
=
\mathrm{Carr}_p\!\Bigl(\bigl(d_{m+k}(u)\bigr)_{u\in re_{m+k\to m}^{-1}(w)}\Bigr).
\]
对 \(w\) 求和即得所示 \(p\)-进重整化公式。证毕。
\end{proof}
\leanverified{paper\_conclusion\_foldbin\_deep\_volume\_carry\_renormalization}

\begin{corollary}[深层 \texorpdfstring{$p$}{p}--进平直性的充要判别]\label{cor:conclusion-foldbin-deep-padic-flatness-criterion}
\leanverified{paper\_conclusion\_foldbin\_deep\_padic\_flatness\_criterion}
沿定理 \ref{thm:conclusion-foldbin-deep-volume-carry-renormalization} 的记号，
\[
v_p\!\left(\left|G_m^{(2^k)}\right|\right)=v_p\!\left(\left|G_{m+k}\right|\right)
\]
当且仅当对每个 \(w\in X_m\)，其全部 \(k\) 层后裔多重度束
\[
\bigl(d_{m+k}(u)\bigr)_{u\in re_{m+k\to m}^{-1}(w)}
\]
在 \(p\)-进制求和时全程无进位。特别地，\(p=2\) 时，跨 \(k\) 层的二进曲率缺陷消失，当且仅当每个祖先纤维的全后裔质量分裂都是逐位无进位的。
\end{corollary}
\begin{proof}
定理 \ref{thm:conclusion-foldbin-deep-volume-carry-renormalization} 右端是若干非负整数之和。该和为零，当且仅当每一项都为零；而 \(\mathrm{Carr}_p=0\) 与逐位无进位完全等价。证毕。
\end{proof}

\begin{proposition}[二值秩一的 restriction--加法缺陷]\label{prop:conclusion-foldbin-restriction-addition-rankone-defect}
\leanverified{paper\_conclusion\_foldbin\_restriction\_addition\_rankone\_defect}
在稳定加法
\[
x\oplus_m y:=s_m\!\bigl(V_m(x)+V_m(y)\bmod F_{m+2}\bigr)
\]
下，定义 restriction 缺陷
\[
\partial_{m+1\to m}(x,y)
:=
V_m\!\bigl(\pi_{m+1\to m}(x\oplus_{m+1}y)\bigr)
-
V_m\!\bigl(\pi_{m+1\to m}(x)\oplus_m\pi_{m+1\to m}(y)\bigr)
\in \ZZ/(F_{m+2}\ZZ).
\]
则
\[
\partial_{m+1\to m}(x,y)\in \{0,F_m\},
\qquad
\partial_{m+1\to m}(x,y)=\kappa_{m+1}^{\mathrm{car}}(x,y)\,F_m.
\]
等价地，在稳定类型坐标中，所有 restriction--加法失真都由唯一缺陷元
\[
\chi_m^{\mathrm{car}}=s_m(F_m)
\]
生成，且其像集恰为 \(\{s_m(0),\chi_m^{\mathrm{car}}\}\)。
\end{proposition}
\begin{proof}
正文定理 9.723 已给出严格恒等式
\[
\pi_{m+1\to m}(x\oplus_{m+1}y)
=
\pi_{m+1\to m}(x)\oplus_m\pi_{m+1\to m}(y)\oplus_m
\bigl(\kappa_{m+1}^{\mathrm{car}}(x,y)\cdot \chi_m^{\mathrm{car}}\bigr),
\]
其中 \(\kappa_{m+1}^{\mathrm{car}}(x,y)\in\{0,1\}\)，且 \(\chi_m^{\mathrm{car}}=s_m(F_m)\)。将上式经 \(V_m\) 送入 \(\ZZ/(F_{m+2}\ZZ)\) 即得
\[
\partial_{m+1\to m}(x,y)=\kappa_{m+1}^{\mathrm{car}}(x,y)\,F_m.
\]
故 restriction 失败不是高秩畸变，而是一个单生成、二值触发的缺陷机制。证毕。
\end{proof}
